\documentclass[11pt,a4paper]{article}
\usepackage{operationsform}
\begin{document}
\OPStitle{アクセシビリティ品質チェックリスト}{TPL-CHK-010}{0.2.0}
\OPSsection{1}{確認項目}
\OPScheck{1}{情報を運ぶ画像に目的に応じた代替テキストが存在する}
\OPScheck{2}{複雑なグラフに主要な数値・傾向を説明する代替経路がある}
\OPScheck{3}{音声・動画に必要な字幕・書き起こしと説明文を用意した}
\OPScheck{4}{見出し構造と文書の読み順が内容の関係を表している}
\OPScheck{5}{色、音、動きの一つの感覚手段だけに重要情報を依存させていない}
\OPScheck{6}{専門語・略語の説明が対象読者に適切な位置にある}
\OPScheck{7}{異なる端末・表示環境で読解不能な部分がないか確認する}
\OPScheck{8}{確認者、参照した技術基準、未解決の課題を記録した}
\OPSsection{2}{点検結果の記録}
\OPSblankshort{対象案件・対象版}
\OPSblankshort{担当者と点検日}
\OPSblanklong{未解決事項と対応}
\par\Needspace{4\baselineskip}\noindent\textbf{確認に用いた資料ID}\par\vspace{1mm}
{\fontsize{8.5}{13}\selectfont
\begin{tabularx}{\linewidth}{@{}XXXX@{}}
\textbf{対象・項目} & \textbf{内容・参照} & \textbf{担当・状態} & \textbf{備考}\\[2mm]
\strut & \strut & \strut & \strut\\[5mm]
\strut & \strut & \strut & \strut\\[5mm]
\end{tabularx}
}\vspace{2mm}\par
\end{document}
